\section{Results}
\label{sec:results}

\subsection{Primary Hypothesis Tests}
\label{sec:primary}

\Tabref{tab:hypotheses} presents results for all four hypothesis tests.
Three of four tests strongly refute the artificial outsider hypothesis, while the fourth provides only weak support with a small effect size.

\begin{table}[t]
    \centering
    \caption{Primary hypothesis test results. All tests use Bonferroni-corrected $\alpha = 0.0125$. Three of four tests refute the artificial outsider hypothesis---LLM consensus positively predicts human interestingness. H2 provides weak support (LLM disagreement correlates with Engagement), but the effect size is small.}
    \begin{tabular}{@{}llccl@{}}
        \toprule
        \textbf{Hypothesis} & \textbf{Statistic} & \textbf{Value} & \textbf{$p$-value} & \textbf{Verdict} \\
        \midrule
        H1: Consensus vs.\ Surprise & Spearman $\rho$ & $+0.402$ & $3.1 \times 10^{-42}$ & \textbf{Refuted} \\
        H2: Disagreement vs.\ Engagement & Spearman $\rho$ & $+0.143$ & $3.2 \times 10^{-6}$ & Weak support \\
        H3: Bottom vs.\ Top quartile & Mann-Whitney $U$ & --- & $1.00$ & \textbf{Refuted} \\
        H4: LLM-only consensus vs.\ interest & Spearman $\rho$ & $+0.358$ & $1.8 \times 10^{-30}$ & \textbf{Refuted} \\
        \bottomrule
    \end{tabular}
    \label{tab:hypotheses}
\end{table}

\para{H1: LLM consensus positively predicts human surprise.}
LLM consensus correlates positively with human Surprise ratings ($\rho = +0.402$, $p = 3.1 \times 10^{-42}$), the opposite direction predicted by the hypothesis.
Stories that LLMs collectively rate as interesting are also the ones humans find more surprising.

\para{H2: LLM disagreement weakly predicts engagement.}
LLM disagreement shows a small positive correlation with human Engagement ($\rho = +0.143$, $p = 3.2 \times 10^{-6}$).
While statistically significant, the effect size is small and may reflect the inherent complexity of variable stories rather than a genuine ``outsider'' effect.

\para{H3: Low-consensus stories are less interesting to humans.}
The Mann-Whitney $U$ test yields $p = 1.00$, indicating that stories in the bottom LLM consensus quartile are not rated higher by humans than stories in the top quartile.
In fact, the relationship is reversed, as the quartile analysis below shows.

\para{H4: The effect holds for LLM-generated stories specifically.}
Restricting to LLM-generated stories only, the positive correlation persists ($\rho = +0.358$, $p = 1.8 \times 10^{-30}$), confirming that the finding is not driven by human-written stories.

\subsection{LLM Consensus Quartile Analysis}
\label{sec:quartiles}

\Tabref{tab:quartiles} presents human interestingness ratings broken down by LLM consensus quartiles.
The relationship is monotonically increasing: stories in the top quartile (Q4) receive 47\% higher human interestingness ratings than stories in the bottom quartile (Q1).

\begin{table}[t]
    \centering
    \caption{Human interestingness ratings by LLM consensus quartile. Q1 contains stories rated lowest by LLM judges; Q4 contains stories rated highest. Human ratings increase monotonically with LLM consensus---the opposite of the outsider prediction.}
    \begin{tabular}{@{}lcccc@{}}
        \toprule
        \textbf{LLM Quartile} & \textbf{$N$} & \textbf{Interest (mean $\pm$ SEM)} & \textbf{Surprise} & \textbf{Engagement} \\
        \midrule
        Q1 (low consensus) & 467 & $2.093 \pm 0.026$ & 1.858 & 2.328 \\
        Q2 & 185 & $2.344 \pm 0.040$ & 2.050 & 2.638 \\
        Q3 & 211 & $2.470 \pm 0.039$ & 2.148 & 2.791 \\
        Q4 (high consensus) & 193 & $\mathbf{3.073 \pm 0.053}$ & \textbf{2.720} & \textbf{3.427} \\
        \bottomrule
    \end{tabular}
    \label{tab:quartiles}
\end{table}

\subsection{Quality-Controlled Analysis}
\label{sec:quality}

\Tabref{tab:quality} reveals the critical finding: the strong raw correlation between LLM consensus and human interestingness is almost entirely explained by text quality.

\begin{table}[t]
    \centering
    \caption{Correlation between LLM consensus and human interestingness under different controls. The raw correlation ($\rho = +0.482$) drops to near zero after controlling for human-rated quality, indicating LLMs detect quality rather than interestingness per se.}
    \begin{tabular}{@{}lccl@{}}
        \toprule
        \textbf{Analysis} & \textbf{$\rho$} & \textbf{$p$-value} & \textbf{Interpretation} \\
        \midrule
        Raw correlation & $+0.482$ & $1.7 \times 10^{-62}$ & Strong positive \\
        Partial (controlling quality) & $+0.002$ & $0.943$ & \textbf{Near zero} \\
        Within-model average & $+0.205$ & --- & Weak positive \\
        LLM-only, quality-controlled & $+0.023$ & $0.470$ & Near zero \\
        Surprise, quality-controlled & $+0.010$ & $0.746$ & Near zero \\
        \bottomrule
    \end{tabular}
    \label{tab:quality}
\end{table}

After partialing out human-rated quality (mean of all 6 dimensions), the correlation between LLM consensus and human interestingness drops from $\rho = +0.482$ to $\rho = +0.002$ ($p = 0.943$).
This pattern holds when restricting to LLM-generated stories only ($\rho = +0.023$, $p = 0.470$) and when examining Surprise alone ($\rho = +0.010$, $p = 0.746$).
LLMs are not biased for or against interesting content---they are detecting the same underlying quality signal that humans detect.

\subsection{Extreme Group Comparison}
\label{sec:extreme}

Stories where all three LLM judges rate $\leq 2$ ($N = 463$) receive a mean human interestingness rating of 2.091, while stories where all three rate $\geq 4$ ($N = 66$) receive a mean of 3.548.
Unanimously LLM-disliked stories are not ``hidden gems''---they are genuinely less interesting to humans as well.

\subsection{Inter-LLM Agreement}
\label{sec:agreement}

The three LLM judges show high pairwise agreement: $\rho = 0.768$ (\gptfourmini vs.\ \gptfo), $\rho = 0.766$ (\gptfo vs.\ \gptfomini), and $\rho = 0.707$ (\gptfourmini vs.\ \gptfomini).
This strong inter-judge agreement ($\rho = 0.71$--$0.77$) is consistent with prior work documenting LLM evaluation convergence~\citep{echoes2024} and validates our use of multi-judge consensus as a meaningful signal.

\subsection{Arena Complementary Analysis}
\label{sec:arena}

\Tabref{tab:arena} summarizes results from 200 Chatbot Arena battles.
LLM judges align strongly with human preferences, with 77--84\% agreement depending on the model.
Humans chose the LLM-preferred response in 77.2\% of battles, significantly above chance ($p = 1.7 \times 10^{-11}$, binomial test).
The human-preferred response received a higher mean LLM rating (3.29 vs.\ 2.88; paired $t = 7.88$, $p = 2.1 \times 10^{-13}$).

\begin{table}[t]
    \centering
    \caption{Arena complementary analysis (200 battles). LLM judges align strongly with human pairwise preferences across conversational data.}
    \begin{tabular}{@{}lc@{}}
        \toprule
        \textbf{Metric} & \textbf{Value} \\
        \midrule
        LLM-human agreement rate & 77--84\% \\
        Human-preferred response (LLM mean) & 3.29 \\
        Human-rejected response (LLM mean) & 2.88 \\
        Paired $t$-test & $t = 7.88$, $p = 2.1 \times 10^{-13}$ \\
        Fraction humans chose lower-LLM-rated & 22.8\% \\
        Binomial test vs.\ 50\% & $p = 1.7 \times 10^{-11}$ \\
        \bottomrule
    \end{tabular}
    \label{tab:arena}
\end{table}
